\documentclass[11pt]{article}
\usepackage[margin=1in]{geometry}
\usepackage{amsmath,amssymb,amsthm,graphicx}
\usepackage[hidelinks]{hyperref}
\newtheorem{theorem}{Theorem}
\newcommand{\F}{\mathbb F}
\newcommand{\T}{\mathbb T}
\newcommand{\dI}{d_{\mathrm{isl}}}
\setlength{\parindent}{0pt}
\setlength{\parskip}{0.65em}

\title{A Sharp Perturbation Dichotomy for\\
Independent-Island Phase Retrieval}
\author{Run \texttt{fa\_banach\_001}, lane 11}
\date{22 July 2026}

\begin{document}
\maketitle

\textbf{Status.} Candidate sharp dichotomy (positive theorem plus
counterexample), likely valid, for human review.

\section{Source question}

Alharbi, Freeman, Ghoreishi, Lois, and Sebastian prove that ordinary stable
phase retrieval in a finite-dimensional Hilbert space persists under a small
perturbation of the frame vectors. They then discuss signals supported on
separated islands, where recovery only up to an independent phase on each
island is sufficient, and state that the effect of perturbations in this
setting remains an important open problem
\cite[Introduction, PDF p.~3]{AFGLS}.

\begin{center}
\includegraphics[width=0.65\textwidth]{figures/open_problem_crop.png}
\end{center}

The question is a research program rather than a single quantified
conjecture. We formalize its basic orthogonal-island model and obtain a sharp
dichotomy. Stable recovery need not persist under arbitrary perturbations, but
it does persist---with the same explicit dimension-free estimate as in the
source theorem---when every measurement vector remains in its original
island.

\section{Claimed contribution}

Let $\F$ be either $\mathbb R$ or $\mathbb C$, and put $H=\F^2$ with the
orthogonal island decomposition
\[
 H=H_1\oplus H_2,\qquad H_1=\F e_1,\quad H_2=\F e_2.
\]
Define $x\sim_{\mathrm{isl}}y$ when there are $\lambda_1,\lambda_2\in\T$
(with $\T=\{-1,1\}$ over $\mathbb R$) such that
\[
 x_k=\lambda_k y_k\qquad(k=1,2).
\]
Thus relative phase between the two islands is deliberately ignored. Equip
the quotient with
\[
 \dI(x,y)=\min_{\lambda_1,\lambda_2\in\T}
 \left(|x_1-\lambda_1y_1|^2+|x_2-\lambda_2y_2|^2\right)^{1/2}.
\]

\begin{theorem}\label{thm:main}
The frame $X=(e_1,e_2)$ does $1$-stable phase retrieval modulo
$\sim_{\mathrm{isl}}$. For every $\eta>0$, however, there is a frame
$Y=(y_1,y_2)$ satisfying
\[
 \sum_{j=1}^2\|y_j-e_j\|^2<\eta
\]
which is not even injective modulo $\sim_{\mathrm{isl}}$. Consequently no
analogue of arbitrary-perturbation stability is possible for a fixed
independent-island quotient, even in dimension two.
\end{theorem}

\section{Proof intuition}

The coordinate frame measures the two islands separately, so its magnitude
map is exactly the quotient-coordinate map $(x_1,x_2)\mapsto(|x_1|,|x_2|)$.
Mix a tiny amount of the second island into the first measurement. Flipping
the sign of that mixed measurement and solving the resulting two linear
equations produces another signal with identical magnitudes. The new signal
has a nonzero first-island component, whereas the original one does not, so
the two signals are not equivalent under independent island phases. The
instability is therefore qualitative, not merely a loss in the Lipschitz
constant.

Conversely, if measurement vectors are forbidden from crossing between
islands, the problem separates exactly. Restricting the original stability
inequality to one island gives ordinary stable phase retrieval there. The
source perturbation theorem applies on every island, and orthogonality lets us
sum the squared estimates without losing either the threshold or the
stability constant.

\section{Counterexample proof}

For scalars $a,b$,
\[
 \min_{\lambda\in\T}|a-\lambda b|=\big||a|-|b|\big|.
\]
It follows coordinatewise that
\begin{equation}\label{eq:quotient}
 \dI(x,y)^2=\big(|x_1|-|y_1|\big)^2+
                    \big(|x_2|-|y_2|\big)^2.
\end{equation}
The analysis map of $X=(e_1,e_2)$ is $\Theta_Xx=(x_1,x_2)$, hence
\[
 \dI(x,y)=\big\||\Theta_Xx|-|\Theta_Xy|\big\|_2.
\]
This proves exact $1$-stability modulo $\sim_{\mathrm{isl}}$.

Fix $\eta>0$ and choose $0<\varepsilon<\sqrt\eta$. Set
\[
 y_1=e_1+\varepsilon e_2,\qquad y_2=e_2.
\]
Then
\[
 \|y_1-e_1\|^2+\|y_2-e_2\|^2=\varepsilon^2<\eta.
\]
The analysis matrix of $Y=(y_1,y_2)$ is
\[
 S_\varepsilon=
 \begin{pmatrix}1&\varepsilon\\0&1\end{pmatrix},
\]
so $Y$ remains a basis (and therefore a frame) for every $\varepsilon$.

Now take
\[
 x=e_2=(0,1),\qquad z=-2\varepsilon e_1+e_2=(-2\varepsilon,1).
\]
Both vectors are real, so the same calculation works over either scalar
field. Directly,
\[
 S_\varepsilon x=(\varepsilon,1),
 \qquad S_\varepsilon z=(-\varepsilon,1).
\]
Thus
\[
 |\Theta_Yx|=|\Theta_Yz|=(\varepsilon,1).
\]
On the other hand, equation~\eqref{eq:quotient} gives
\[
 \dI(x,z)=
 \left((0-2\varepsilon)^2+(1-1)^2\right)^{1/2}
 =2\varepsilon>0.
\]
Therefore $x\not\sim_{\mathrm{isl}}z$ despite identical magnitude
measurements. The perturbed frame fails injectivity on the quotient, so it
cannot satisfy any finite stability inequality. This proves the theorem.

\section{Structure-preserving perturbations}

We now prove the positive half in the natural finite-dimensional
orthogonal-island model. Let
\[
 H=\bigoplus_{k=1}^r H_k
\]
be an orthogonal decomposition, let $P_k$ denote the orthogonal projection
onto $H_k$, and partition the frame index set as
$J=J_1\sqcup\cdots\sqcup J_r$. Define the independent-island quotient metric
by
\begin{equation}\label{eq:blockmetric}
 d_r(x,y)=\min_{\lambda_1,\ldots,\lambda_r\in\T}
 \left\|x-\sum_{k=1}^r\lambda_kP_ky\right\|.
\end{equation}

\begin{theorem}[Island-preserving perturbation theorem]\label{thm:positive}
Let $X=(x_j)_{j\in J}$ be a frame for $H$ with frame bounds
$B\geq A>0$. Assume that
\[
 x_j\in H_k\quad(j\in J_k)
\]
and that $X$ does $C$-stable phase retrieval for the quotient metric $d_r$:
\[
 d_r(x,y)\leq C
 \left(\sum_{j\in J}
 \big||\langle x,x_j\rangle|-|\langle y,x_j\rangle|\big|^2\right)^{1/2}.
\]
Let
\[
 0<\nu<\frac{1}{16C^4B},
\]
and let $Y=(y_j)_{j\in J}$ be a structure-preserving perturbation satisfying
\[
 y_j\in H_k\quad(j\in J_k),
 \qquad \sum_{j\in J}\|x_j-y_j\|^2<\nu.
\]
Then $Y$ is a frame with bounds
\[
 A\left(1-\sqrt{\frac{\nu}{A}}\right)^2
 \quad\text{and}\quad
 B\left(1+\sqrt{\frac{\nu}{B}}\right)^2,
\]
and it does $C_\nu$-stable phase retrieval modulo independent island phases,
where
\[
 C_\nu=C\left(1-4C^2\sqrt{\nu B}\right)^{-1/2}.
\]
Thus the threshold and stability estimate are exactly those of the ordinary
perturbation theorem in the source.
\end{theorem}

\begin{proof}
For $u,v\in H_k$, equation~\eqref{eq:blockmetric} reduces to the ordinary
phase quotient:
\[
 d_r(u,v)=\min_{\lambda\in\T}\|u-\lambda v\|.
\]
Measurements indexed outside $J_k$ vanish on $H_k$. Hence the restriction
$X_k=(x_j)_{j\in J_k}$ is a frame for $H_k$ with the same admissible frame
bounds $A,B$, and the assumed quotient inequality shows that $X_k$ does
$C$-stable ordinary phase retrieval.

Put $Y_k=(y_j)_{j\in J_k}$. Its perturbation size satisfies
\[
 \sum_{j\in J_k}\|x_j-y_j\|^2<\nu.
\]
The ordinary perturbation theorem of Alharbi et al. therefore applies to
every pair $X_k,Y_k$ with the same parameter $\nu$. It gives the displayed
frame bounds and, for all $u,v\in H_k$,
\begin{equation}\label{eq:blockestimate}
 \min_{\lambda\in\T}\|u-\lambda v\|
 \leq C_\nu\left(\sum_{j\in J_k}
 \big||\langle u,y_j\rangle|-|\langle v,y_j\rangle|\big|^2\right)^{1/2}.
\end{equation}

Summing the block frame inequalities proves the asserted global frame bounds
for $Y$. Finally, orthogonality makes the minimization in
\eqref{eq:blockmetric} separate:
\[
 d_r(x,y)^2=\sum_{k=1}^r
 \min_{\lambda\in\T}\|P_kx-\lambda P_ky\|^2.
\]
Apply the square of \eqref{eq:blockestimate} with
$u=P_kx$, $v=P_ky$, and sum over $k$. Since $y_j\in H_k$ for $j\in J_k$,
the right-hand side becomes
\[
 C_\nu^2\sum_{j\in J}
 \big||\langle x,y_j\rangle|-|\langle y,y_j\rangle|\big|^2.
\]
Taking square roots gives the claimed global stability estimate.
\end{proof}

\section{Sharp dichotomy and scope}

The example embeds in $\F^m$ for every $m\ge2$: retain the same first two
frame vectors and signals, and append $e_3,\dots,e_m$. Hence neither dimension
nor the number of separated islands repairs the obstruction.

Theorem~\ref{thm:main} and Theorem~\ref{thm:positive} give the promised sharp
answer in the orthogonal-island model. Allowing even one arbitrarily small
cross-island component can destroy quotient injectivity, whereas every
sufficiently small island-preserving perturbation retains stable recovery with
an explicit dimension-free constant. More general notions of approximately
separated islands, signal-dependent island decompositions, or perturbations
that change the quotient itself remain separate questions.

\section{Exact verification}

The script \texttt{code/verify\_counterexample.py} checks with exact rational
arithmetic that the perturbation size is nonzero and arbitrarily scalable,
the perturbed vectors form a basis, the two displayed magnitude measurements
agree, and their quotient distance is positive. The positive theorem is an
exact blockwise consequence of the source perturbation theorem; its only
additional identity is the orthogonal sum-of-squares decomposition displayed
in the proof.

\begin{thebibliography}{9}
\bibitem{AFGLS}
W.~Alharbi, D.~Freeman, D.~Ghoreishi, C.~Lois, and S.~Sebastian,
\emph{Stable phase retrieval and perturbations of frames},
Proc. Amer. Math. Soc. Ser. B 10 (2023), 353--368;
\href{https://arxiv.org/abs/2212.13681}{arXiv:2212.13681}.
\end{thebibliography}

\end{document}
